\section{Proof of strict positivity}\label{sec:positivity}

We apply the two norm comparisons to different finite extensions.  The
first comparison replaces one of the original modules by a normalized
length algebra.  The second comparison is applied to bounded complexes
obtained from finite approximations to a parameter quotient of that
algebra.  The normalized lengths need not be invariant under either
Galois group.

\subsection{Reduction and the first norm comparison}

\begin{lemma}\label{pos:reduction}
To prove Theorem~\ref{thm:main}, it suffices to prove the assertion for
\[
 D=R/P,\qquad E=R/Q,\qquad
 h=\dim D>0,\qquad l=\dim E>0,\qquad h+l=d,
\]
where $R$ is complete with algebraically closed residue field $k$ of
characteristic $p>0$, $P,Q$ are prime, and $P+Q$ is primary to $\m_R$.
\end{lemma}

\begin{proof}
The equicharacteristic case is Serre's positivity theorem
\cite{Serre}.  We may therefore assume that the residue characteristic
is positive.  Completion, followed by the extension in
Lemma~\ref{len:residue-extension}, reduces to a complete regular local
ring with algebraically closed residue field.  For clarity, a flat
local extension $R\to R_1$ with $\m_RR_1=\m_{R_1}$ preserves the
length of every finite-length module: a composition series becomes a
filtration with factors $R_1/\m_{R_1}$.  Flat base change identifies the
Tor modules after extension.  The same argument applied to
$M/\m_R^nM$ preserves the Hilbert--Samuel function, and hence the
dimension of $M$; it applies equally to $N$.  It also preserves
regularity, since it preserves the dimension and the minimal number
of generators of the maximal ideal.

We use the vanishing theorem of Roberts and Gillet--Soul\'e in its
general regular local form \cite{Roberts,GilletSoule}:
\begin{equation}\label{pos:vanishing}
 \chi^R(V,W)=0
 \quad\text{if}\quad
 \length_R(V\otimes_R W)<\infty
 \quad\text{and}\quad \dim V+\dim W<d.
\end{equation}
Take prime filtrations of $M$ and $N$.  All pairs of factors still
have finite-length tensor product, since their supports are contained
in the respective original supports.  Additivity of the Euler
characteristic and Equation~\eqref{pos:vanishing} discard every pair except
those having both original dimensions.  There are top-dimensional
factors in each filtration, so at least one such pair remains.
Proving strict positivity for each remaining pair proves it for
$M,N$.  Their primes satisfy $\sqrt{P+Q}=\m_R$.

Finally, a zero-dimensional prime quotient is $k$, and its
complementary $d$-dimensional prime quotient is $R$, since a regular
local ring is a domain.  For that pair the characteristic equals
$\length_R(k)=1$.  Thus we may assume $h,l>0$.
\end{proof}

Fix a pair as in Lemma~\ref{pos:reduction}.  By
Lemma~\ref{len:normalization}, choose regular complete normalization
bases
\[
 A_D\subset D,\qquad A_E\subset E
\]
of dimensions $h,l$, with residue field $k$, such that the extensions
are finite and generically separable.

Let $D'$ and $E'$ be the integral closures of these bases in finite
Galois closures of the respective fraction-field extensions.  They
are finite over their bases by excellence, and are finite over $D,E$
as well.  They are local: a finite algebra over a complete local ring
is a product of complete local algebras, and a domain has only one
factor \cite[Tag~03QH]{Stacks}. Their residue fields are finite extensions of $k$, hence are
$k$.  Write
\[
 G_D=\operatorname{Gal}(\operatorname{Frac}(D')/
                  \operatorname{Frac}(A_D)),\qquad
 G_E=\operatorname{Gal}(\operatorname{Frac}(E')/
                  \operatorname{Frac}(A_E)).
\]
Normality of the bases gives $(D')^{G_D}=A_D$ and
$(E')^{G_E}=A_E$. For either extension write $T^G=A$.
Let $\Omega$ be an algebraically closed $A$-field and put
$S=T\otimes_A\Omega$. For $u=\sum_i t_i\otimes\omega_i\in S$,
the norm $\prod_{g\in G}g(u)$ is scalar: the coefficients of
$\prod_{g\in G}(\sum_i g(t_i)X_i)$ lie in $T^G=A$, and we
specialize $X_i=\omega_i$. If the finite Artinian $\Omega$-algebra
$S$ had more than one orbit of points, its product decomposition
would give a nontrivial $G$-invariant idempotent $e$, equal to one on
one orbit and zero on the others. But its norm is $e^{|G|}=e$ and is
scalar, so $e=0$ or $e=1$, a contradiction. Thus every geometric fiber
is transitive, without any assertion that invariants commute with
base change. The prime and residue-field formulation is also given
in \cite[Tags~0BRI and 0BRJ]{Stacks}. Thus the hypotheses of
Theorem~\ref{ds:descent} apply to both extensions.

Use Theorem~\ref{len:input} to choose the algebras and normalized
lengths
\[
 (C_D,\lambda_D)=(C_{D'},\lambda_{D'}),\qquad
 (C_E,\lambda_E)=(C_{E'},\lambda_{E'}).
\]
They remain fixed throughout the proof.  Set
$a_D=\rank_D D'$ and $a_E=\rank_E E'$.  For $\sigma\in G_D$,
define
\[
 \iota_\sigma:R\longrightarrow D'
      \xrightarrow{\sigma}D'.
\]
Give $D'$ this $R$-action, and give $C_D$ the action obtained by
composing $\iota_\sigma$ with $D'\to C_D$. Define
\begin{equation}\label{pos:chi-sigma}
 c_\sigma
   =\sum_{i=0}^d(-1)^i
      \lambda_D\bigl(\Tor_i^R(C_D,E')\bigr),
\end{equation}
where this $R$-action on $C_D$ is understood.

\begin{lemma}\label{pos:first-average}
The lengths in Equation~\eqref{pos:chi-sigma} are finite, and
\begin{equation}\label{pos:first-average-equation}
 |G_D|a_Da_E\,\chi^R(D,E)=\sum_{\sigma\in G_D}c_\sigma.
\end{equation}
\end{lemma}

\begin{proof}
Take a finite $R$-free resolution of $E'$ and extend it to $D'$.
The resulting bounded free complex is supported at the closed point:
after localizing its homology away from $\m_R$, the supports of
$D'$ and $E'$ are disjoint.  The same holds after every twist, whose
kernel on $R$ is still $P$.  Lemma~\ref{ds:ordinary-norm} supplies
finite normalized homology lengths and equality of the ordinary and
normalized Euler evaluations after summing all conjugates.

Applying $\sigma$ entrywise gives a semilinear isomorphism from the
untwisted complex to its conjugate, so all ordinary Euler
characteristics in that sum are equal to $\chi^R(D',E')$.  Ordinary
lengths over $D'$ and over $R$ agree here because both residue fields
are $k$.

Choose an injection $D^{a_D}\to D'$ with torsion cokernel.  The
cokernel has dimension less than $h$, and therefore its intersection
characteristic with $E'$ vanishes by Equation~\eqref{pos:vanishing}.
Additivity gives
\[
 \chi^R(D',E')=a_D\chi^R(D,E').
\]
The analogous argument for $E\subset E'$ gives
$\chi^R(D,E')=a_E\chi^R(D,E)$.  The norm comparison now yields
Equation~\eqref{pos:first-average-equation}.
\end{proof}

It remains to show $c_\sigma>0$ for each $\sigma$.  Fix one
$\sigma$ until the final paragraph, and write $C=C_D$ and
$\lambda=\lambda_D$ with the chosen $R$-action.

\subsection{A parameter quotient and its Tor groups}

Choose
\begin{equation}\label{pos:parameters}
 x_1,\ldots,x_h\in Q
\end{equation}
which form an $R$-regular sequence and whose images are a system of
parameters of $D$.  Here is the simultaneous prime-avoidance
argument.  After $i-1$ choices, the associated primes of
$R/(x_1,\ldots,x_{i-1})$ have height $i-1$, since $R$ is
Cohen--Macaulay.  None contains $Q$, whose height is $h\ge i$.
Also avoid the primes in $D/(x_1,\ldots,x_{i-1})$ which obstruct the
next parameter dimension drop.  Each has positive dimension; if its
inverse image contained $Q$, it would contain $P+Q$ and hence would
be the maximal ideal, a contradiction.  Prime avoidance inside $Q$
therefore makes both choices possible at every step.

Put
\[
 B=R/(x_1,\ldots,x_h),\qquad
 I_\sigma=(\iota_\sigma(x_1),\ldots,
                 \iota_\sigma(x_h))D',\qquad U=C/I_\sigma C.
\]
The ideal $I_\sigma$ is a parameter ideal of $D'$, since finiteness
and then the automorphism $\sigma$ preserve that property.  Thus
Theorem~\ref{len:input} gives
\begin{equation}\label{pos:U-positive}
 0<\lambda(U)=e(I_\sigma,D')<\infty,
 \qquad
 \lambda\bigl(H_q(x;C)\bigr)=0\quad(q>0).
\end{equation}
Here the notation $x$ in a complex over $C$ uses $\iota_\sigma$.
Some integer $a\ge1$ satisfies
\begin{equation}\label{pos:uniform-annihilator}
 \m_B^a(D'/I_\sigma)=0.
\end{equation}
In particular it annihilates $U$ and every module over $D'/I_\sigma$.

\begin{lemma}\label{pos:tor-bound}
For every finite $B$-module $L$, all the following lengths are finite,
and
\begin{equation}\label{pos:tor-identification}
 \tau_i(L):=\lambda\bigl(\Tor_i^B(U,L)\bigr)
     =\lambda\bigl(\Tor_i^R(C,L)\bigr),
 \qquad \tau_i(L)=0\quad(i>l).
\end{equation}
\end{lemma}

\begin{proof}
The complexes in question have compatible $C$-actions.  The Koszul
homology $H_q(x;C)$ is annihilated by $x$ and by $P$, and is therefore
bounded-primary-torsion.  For $q>0$ its length is zero by
Equation~\eqref{pos:U-positive}.  A free $B$-resolution of $L$,
with each term of finite rank, gives
\[
 \lambda\bigl(\Tor_p^B(H_q(x;C),L)\bigr)=0\quad(q>0).
\]
Indeed every such homology is a subquotient of a finite direct sum
of a length-zero module.  For $q=0$ these lengths are finite because
$\lambda(U)<\infty$.  The change-of-rings spectral sequence
\[
 \Tor_p^B\bigl(\Tor_q^R(C,B),L\bigr)
       \ \Longrightarrow\ \Tor_{p+q}^R(C,L)
\]
now gives the asserted equality of lengths.  Only finitely many
terms enter each total degree.

For the sharper bound, choose an $R$-free resolution $P_L$ of $L$
in degrees $0,\ldots,d$ and consider
\[
 K(x;R)\otimes_R C\otimes_R P_L.
\]
Taking Koszul homology first, all positive rows have length zero,
and the zero row is $U\otimes_R P_L$.  Consequently the total
homology has length zero in degrees greater than $d$.  On the other
hand, $xL=0$, so in the derived category of $R$-modules
\[
 K(x;R)\otimes_R^{\mathbf L}L
       \simeq\bigoplus_{j=0}^h L[j]^{\binom hj}.
\]
Tensoring this identity derivedly with $C$ preserves the splitting
as a statement about $C$-complexes.  The total homology in degree
$i+h$ therefore contains $\Tor_i^R(C,L)$ as a direct summand.  If
$i>l=d-h$, that total degree exceeds $d$, proving the bound.  This
argument also shows directly that all the displayed lengths are
finite.
\end{proof}

\subsection{Finite approximations preserving the actions}

We need finite-length $B$-modules whose normalized Tor lengths approach
those of $U$, while retaining the single annihilator
$\m_B^a$.  No finite generation or finite presentation of $U$ is
assumed.  The approximation will take place among modules carrying
both its tower action and its $B$-action.  In this subsection,
$\lambda$ denotes $\lambda_\infty$ on underlying tower modules,
as permitted by Theorem~\ref{len:input}; the approximating modules
need not carry a $C$-action.

Let $A=A_D$, let $A_j$ be its flat regular tower from
Lemma~\ref{len:tower}, and let $\Lambda=\bigcup_j A_j$.
Thus
\[
 s(j):=\rank_A A_j=p^{jh}.
\]
Set
\begin{equation}\label{pos:approximation-algebras}
 S_j=A_j\otimes_A(D'/I_\sigma),\qquad
 S=\Lambda\otimes_A(D'/I_\sigma).
\end{equation}
The commuting $\Lambda$- and $D'/I_\sigma$-actions on $U$ combine
into an $S$-action.  Working with $S$-modules retains both of these
actions.  Descent to $S_j$ retains the existing $B$-action by
restriction along $B\to D'/I_\sigma$ induced by $\iota_\sigma$;
no new $B$-action is chosen at a finite stage.

\begin{lemma}\label{pos:approximation}
There are finite-length $B$-modules $M_n$ and positive integers
$s_n\to\infty$, with $\m_B^aM_n=0$, such that
\begin{align}
 \frac{\length_B(M_n)}{s_n}&\longrightarrow\lambda(U),
       \label{pos:approximation-length}\\
 \frac{\length_B\Tor_i^B(M_n,L)}{s_n}
       &\longrightarrow\tau_i(L)
       \quad(L=k,E',\ i\ge0).
       \label{pos:approximation-tor}
\end{align}
Writing $\beta_0(M_n)=\dim_k(M_n/\m_BM_n)$, we also have
\begin{equation}\label{pos:generator-bound}
 \frac{\beta_0(M_n)}{s_n}\longrightarrow\nu_\sigma,
 \qquad
 \nu_\sigma\ge
 \frac{\lambda(U)}{\length_B(B/\m_B^a)}>0.
\end{equation}
\end{lemma}

\begin{proof}
The algebra $S$ is finitely presented as a $\Lambda$-module and
has finite normalized length.  Every $S_j$ is finite over the
twisted image of $B$: $D'$ is finite over the twisted image of
$R/P$, and $A_j$ is finite free over $A$.  Equation
\eqref{pos:uniform-annihilator} annihilates $S_j$ and $S$.
Moreover, $S_j$ is Artinian local with residue field $k$.  To see
locality, the tower parameter roots are nilpotent over the
Artinian quotient $D'/I_\sigma$, and reduction leaves the common
residue field $k$.  Thus finite $S_j$-modules have finite lengths
over both $B$ and $A_j$, and these lengths agree: a composition
series over $S_j$ has factors $k$, of length one for either ring.
This argument is unchanged if a residue-field identification is
composed with an automorphism.

Given $\epsilon>0$, choose a finitely generated $S$-submodule
$V\subset U$ with
$\lambda(U/V)<\epsilon$.  This follows from the supremum property
of normalized length: the $S$-span of a finitely generated
$\Lambda$-submodule is still finitely generated over $\Lambda$.
To obtain a finite presentation as well, choose a surjection $S^r\to V$
with kernel $K$.  Its kernel has finite length, bounded by
$r\lambda(S)$.  Choose a finitely generated $S$-submodule $K_0\subset K$ such
that $\lambda(K/K_0)<\epsilon$.  Then
\[
 P_\epsilon=S^r/K_0\longrightarrow U
\]
has kernel and cokernel of length less than $\epsilon$.
It is finitely presented over $S$ and over $\Lambda$, and is
still annihilated by $\m_B^a$.

Fix $B$-resolutions of $k$ and $E'$ with finite free terms, and write
$r_i(L)$ for the rank in degree $i$, with $r_{-1}(L)=0$.
If $K_\epsilon$ and $W_\epsilon$ are the kernel and cokernel of
$P_\epsilon\to U$, the two long exact homology sequences give
\begin{align}
 &\left|\lambda\Tor_i^B(P_\epsilon,L)
                -\lambda\Tor_i^B(U,L)\right| \notag\\
 &\quad\le
   \bigl(r_i(L)+r_{i-1}(L)\bigr)\lambda(K_\epsilon)
   +\bigl(r_{i+1}(L)+r_i(L)\bigr)\lambda(W_\epsilon).
       \label{pos:tor-approximation-error}
\end{align}
For example, each homology of the complex obtained by tensoring
the resolution with $K_\epsilon$ has length at most the rank of
the corresponding term times $\lambda(K_\epsilon)$; the same
holds for $W_\epsilon$.  All these quantities are finite.  Also
$|\lambda(P_\epsilon)-\lambda(U)|<2\epsilon$.

A finite $S$-presentation of $P_\epsilon$ descends to an
$S_j$-module $P_j$ for some $j$, with
$P_\epsilon=\Lambda\otimes_{A_j}P_j$.  Flatness of
$\Lambda/A_j$ gives, for every $i$,
\[
 \Tor_i^B(P_\epsilon,L)
       =\Lambda\otimes_{A_j}\Tor_i^B(P_j,L).
\]
The finite-stage homology on the right is finite over $S_j$.
The normalization formula in Theorem~\ref{len:input}, and the
length comparison above, therefore give exactly
\[
 \lambda\Tor_i^B(P_\epsilon,L)
    =\frac{\length_{A_j}\Tor_i^B(P_j,L)}{s(j)}
    =\frac{\length_B\Tor_i^B(P_j,L)}{s(j)}.
\]
The same formula holds for $P_\epsilon$ itself.

At step $n$, choose $\epsilon$ so that all errors in
Equation~\eqref{pos:tor-approximation-error}, for $i\le n$ and
$L=k,E'$, and the length error are less than $1/n$.
Descend this presentation and enlarge its stage if necessary so
that $j\ge n$.  Further flat tower extension leaves every
normalized value unchanged.  Take $M_n=P_j$ and $s_n=s(j)$.
This proves Equations~\eqref{pos:approximation-length} and
\eqref{pos:approximation-tor}, including the common annihilator
and $s_n\to\infty$.

The case $i=0,L=k$ gives the limit in
Equation~\eqref{pos:generator-bound}.  A minimal generating set and the
annihilator give a surjection
$(B/\m_B^a)^{\beta_0(M_n)}\to M_n$.  Consequently
\[
 \length_B(M_n)
    \le\beta_0(M_n)\length_B(B/\m_B^a).
\]
Divide by $s_n$ and take limits to obtain the stated positive
lower bound.
\end{proof}

\subsection{Bounded complexes and their ordinary evaluation}

The approximations provide two kinds of control.  In minimal
resolutions, the Tor bound $\tau_i(k)=0$ for $i>l$ makes the
corresponding ranks negligible.  The common annihilator then makes
the truncated sequences
supported after quotienting by negligible ranks.  That same annihilator
has already given the positive normalized generator bound
\eqref{pos:generator-bound}, which will make the second evaluation
strictly positive.

Choose minimal free $B$-resolutions $F_n^\infty$ of $M_n$, indexed
homologically in nonnegative degrees.  Let $F_n$ be their brutal
truncations to degrees $0,\ldots,l$.  Minimality and
Equation~\eqref{pos:approximation-tor} give, in each fixed degree,
\begin{equation}\label{pos:rank-estimates}
 \rank_B F_{n,i}^\infty=O(s_n),\qquad
 \rank_B F_{n,i}^\infty=o(s_n)\quad(i>l),
\end{equation}
because these ranks equal $\length_B\Tor_i^B(M_n,k)$.

Let $z_1,\ldots,z_l$ be regular parameters of $A_E$.  For
$\gamma\in G_E$, extend $F_n$ through
\[
 B\longrightarrow E\longrightarrow E'
                  \xrightarrow{\gamma}E',
\]
and denote the resulting complex by $F_{n,\gamma}'$.
Use the normalization $s_n$ in the asymptotic category over $E'$.
These numbers come from the $D$-side approximation. Section~\ref{sec:asymptotic}
allows any positive integer sequence $s_n$, so no identification or
compatibility between the parameter-root towers of $A_D$ and $A_E$
is needed.

\begin{lemma}\label{pos:supported-truncation}
The sequence $(F_{n,\gamma}')_n$ is supported at the closed point
in that category.  Its ordinary Euler evaluation is $c_\sigma$
for every $\gamma\in G_E$.
\end{lemma}

\begin{proof}
Every element $b\in\m_B^a$ acts nullhomotopically on
$F_n^\infty$, because it annihilates $M_n$.  Restrict a chosen
homotopy to the brutal truncation.  Its only remaining error is
in degree $l$, where it is
\begin{equation}\label{pos:truncation-error}
 F_{n,l}\xrightarrow{h_l}F_{n,l+1}^\infty
                      \xrightarrow{d_{l+1}}F_{n,l}.
\end{equation}
This is a factorization by chain maps through
$F_{n,l+1}^\infty[l]$: the map induced by $d_{l+1}$ is a chain map
because $d_ld_{l+1}=0$, and the map induced by $h_l$ is a chain map
because the truncation has no degree $l+1$ term.  The middle module has rank
$o(s_n)$ by Equation~\eqref{pos:rank-estimates}, and is therefore killed
in the asymptotic quotient.  These identities persist after
extension to $E'$, without a flatness assumption over $B$.

The ideal $\m_BE'$ for each indicated structural map is primary
to $\m_{E'}$.  Choose $v$ so large that every $z_j^v$ belongs to
the extension of $\m_B^a$ for all $\gamma$.  This is one fixed
choice, since the group is finite.  Expressing $z_j^v$ as a
finite $E'$-linear combination of elements of $\m_B^a$ proves
that it acts by zero on the sequence in the quotient.  The
same combinations give genuine nullhomotopies on the full
extended resolutions.  This proves support.

Apply the ordinary version of
Equation~\eqref{as:koszul-polynomial}.  In its right-hand side, one may
replace the truncated extended resolutions by the full ones in
each fixed homological degree.  Indeed, the quotient of the
full resolution by its brutal truncation has only terms of
degrees at least $l+1$, all of rank $o(s_n)$ in each fixed degree.
After tensoring that tail with $K(z^v;E')$, taking Koszul
homology first bounds each fixed-degree homology by finite sums
of
\[
 \rank_B F_{n,i}^\infty\,
     \length_{E'} H_q(z^v;E')=o(s_n).
\]
The Koszul homology lengths are fixed and finite, and only
finitely many terms contribute to each total degree.  The long
exact homology sequence thus proves the replacement assertion.

On each full extended resolution the powers $z_j^v$ act
nullhomotopically, as established above.  Koszul tensoring
therefore gives the direct sum of shifts with multiplicities
$\binom lj$.  Its underlying homology has ordinary lengths
$\length_B\Tor_i^B(M_n,E')$, unchanged by conjugating the
$E'$-action: applying $\gamma$ entrywise is a semilinear
isomorphism.  Equations~\eqref{pos:approximation-tor} and
\eqref{pos:tor-identification}, followed by coefficient
comparison in Equation~\eqref{as:koszul-polynomial}, now identify the
ordinary homology-length coefficients with
\[
 \tau_i(E')=\lambda\Tor_i^R(C,E').
\]
Their alternating sum is exactly $c_\sigma$.
\end{proof}

\subsection{A positive normalized evaluation}

The ordinary evaluation has recovered $c_\sigma$ on every conjugate.
We now evaluate the same supported object using $C_E$.  The parameter
Koszul identities will leave only its nonnegative degree-zero
coefficient, and minimality of the resolutions will bound that
coefficient below using \eqref{pos:generator-bound}.

Corollary~\ref{len:finite-length-base-change} gives
\begin{equation}\label{pos:residue-positive}
 0<\eta_E:=\lambda_E(C_E/\m_{E'}C_E)<\infty.
\end{equation}
We record a quantitative bound for use below.  Set
$q_E=(z_1,\ldots,z_l)E'$,
$r_E=\rank_{A_E}E'$, and
$t_E=\length_{E'}(E'/q_E)$.  The length input gives
$\lambda_E(C_E/q_EC_E)=r_E>0$.
A composition series of $E'/q_E$, tensored right-exactly with
$C_E$, shows
\begin{equation}\label{pos:residue-bound}
 r_E\le t_E\eta_E.
\end{equation}
At each step the image of the tensor of the submodule has
length at most that tensor's length; injectivity after tensoring
is not required.  Finiteness of $\eta_E$ follows from the
surjection $C_E/q_EC_E\to C_E/\m_{E'}C_E$.

\begin{lemma}\label{pos:positive-evaluation}
The $\lambda_E$ Euler evaluation of $(F_{n,\gamma}')_n$ is a
number $b_{\sigma,\gamma}$ satisfying
\begin{equation}\label{pos:strict-bound}
 b_{\sigma,\gamma}
    \ge\nu_\sigma\eta_E
    \ge
       \frac{\lambda_D(U)}{\length_B(B/\m_B^a)}
       \frac{r_E}{t_E}
    >0.
\end{equation}
\end{lemma}

\begin{proof}
Let $b_i$ be the nonnegative homology-length coefficients of
this evaluation in Equation~\eqref{as:koszul-polynomial}.  Consider the
finite double complex
\[
 K(z^v;E')\otimes_{E'}F_{n,\gamma}'
                         \otimes_{E'}C_E.
\]
Taking Koszul homology first, Theorem~\ref{len:input} makes all
positive Koszul rows have length zero.  The remaining row has
complex degrees $0,\ldots,l$.  Thus the normalized homology
lengths of this double complex vanish in total degrees greater
than $l$, as well as in negative degrees.  Equation
\eqref{as:koszul-polynomial} says that their polynomial is
$(1+t)^l\sum_i b_it^i$.  All $b_i$ are nonnegative.  Boundedness
and comparison of the least degree rule out negative-degree
coefficients.  Any $b_i>0$ with $i>0$ would then give a positive
coefficient in degree $i+l>l$, a contradiction.  Consequently
only $b_0$ can be nonzero, and the Euler evaluation equals $b_0$.

In degree zero, the displayed double complex has homology
\[
 C_E^{\beta_0(M_n)}\big/
 \left(\operatorname{im}(d_1)+(z_1^v,\ldots,z_l^v)
                                    C_E^{\beta_0(M_n)}\right).
\]
Since $l>0$, $d_1$ is the degree-one presentation differential
of $M_n$.  Minimality over $B$ puts all its entries in $\m_B$,
whose image under any twist is contained in $\m_{E'}$.
The base parameter powers are also in $\m_{E'}$.  There is
therefore a surjection from this degree-zero homology onto
\[
 (C_E/\m_{E'}C_E)^{\beta_0(M_n)}.
\]
Its normalized length is at least
$\beta_0(M_n)\eta_E/s_n$.  The degree-zero coefficient in
Equation~\eqref{as:koszul-polynomial} is $b_0$, so taking the ultralimit
and using the ordinary limit in Equation~\eqref{pos:generator-bound}
gives $b_0\ge\nu_\sigma\eta_E$.  Equations
\eqref{pos:generator-bound} and \eqref{pos:residue-bound} give
the explicit remaining inequalities in
Equation~\eqref{pos:strict-bound}.  Define $b_{\sigma,\gamma}=b_0$.
\end{proof}

\begin{proof}[Proof of Theorem~\ref{thm:main}]
By Lemma~\ref{pos:reduction}, use the domain pair fixed above.
For each $\sigma\in G_D$, carry out the construction with its
specified $R$-action.  Apply Theorem~\ref{as:norm} to the supported
sequence of
Lemma~\ref{pos:supported-truncation}, over $A_E\subset E'$.
It equates the sums of the ordinary and normalized evaluations
of all $G_E$-conjugates.  The former evaluation is $c_\sigma$
for every conjugate, whereas the latter is
$b_{\sigma,\gamma}$.  Hence
\begin{equation}\label{pos:second-average}
 |G_E|c_\sigma
       =\sum_{\gamma\in G_E}b_{\sigma,\gamma}>0
\end{equation}
by Lemma~\ref{pos:positive-evaluation}.
Finally, Lemma~\ref{pos:first-average} gives
\[
 |G_D|a_Da_E\,\chi^R(D,E)
    =\frac{1}{|G_E|}
       \sum_{\sigma\in G_D}\sum_{\gamma\in G_E}
                   b_{\sigma,\gamma}>0.
\]
Both ranks and both group orders are positive.  Thus
$\chi^R(D,E)>0$, and the reduction proves the theorem.
Neither norm comparison requires a Galois action on $C_D$ or
$C_E$, nor invariance of either normalized length under a twist.
\end{proof}
